% =============================================================================
% preamble.tex — gói và cấu hình dùng chung cho report
% Engine: pdflatex. Tiếng Việt qua babel + mã font T5.
% =============================================================================

% --- Lọc warning vô hại (phải nạp trước các gói phát ra warning) ---------------
\usepackage{silence}
\WarningFilter{csquotes}{No space factor codes}          % csquotes chưa biết mã T5
\WarningFilter{microtype}{I cannot find a protrusion list} % font \texttt ở mã T5

% --- Mã hoá và font -----------------------------------------------------------
\usepackage[utf8]{inputenc}
\usepackage[T5]{fontenc}          % T5 có đủ chữ Việt có dấu (T1 thì không)
\usepackage{lmodern}              % Latin Modern: hỗ trợ T5 đầy đủ
\usepackage[vietnamese]{babel}    % Tên tự động: Hình, Bảng, Tài liệu, Chứng minh...

% --- Bố cục trang -------------------------------------------------------------
\usepackage[a4paper, margin=2.5cm, headheight=14pt]{geometry} % headheight theo yêu cầu fancyhdr ở cỡ chữ 12pt
\usepackage{microtype}
\usepackage{setspace}
\setstretch{1.15}
\usepackage{fancyhdr}
\pagestyle{fancy}
\fancyhf{}
\fancyhead[L]{\small\nouppercase{\leftmark}}
\fancyfoot[C]{\small\thepage}
\renewcommand{\headrulewidth}{0.4pt}

% Chống dòng mồ côi/goá (long-form best practices)
\widowpenalty=10000
\clubpenalty=10000

% --- Toán ---------------------------------------------------------------------
\usepackage{mathtools}            % nạp amsmath
\usepackage{amssymb}
\usepackage{amsthm}
\usepackage{bm}

\theoremstyle{plain}
\newtheorem{theorem}{Định lý}[section]
\newtheorem{lemma}[theorem]{Bổ đề}
\newtheorem{proposition}[theorem]{Mệnh đề}
\newtheorem{corollary}[theorem]{Hệ quả}
\theoremstyle{definition}
\newtheorem{definition}[theorem]{Định nghĩa}
\newtheorem{example}[theorem]{Ví dụ}
\theoremstyle{remark}
\newtheorem{remark}[theorem]{Nhận xét}

% Lệnh toán dùng nhiều
\newcommand{\R}{\mathbb{R}}
\newcommand{\N}{\mathbb{N}}
\newcommand{\E}{\mathbb{E}}
\newcommand{\Prob}{\mathbb{P}}
\newcommand{\ind}{\mathbf{1}}
\DeclareMathOperator*{\argmax}{arg\,max}
\DeclareMathOperator*{\argmin}{arg\,min}
\DeclarePairedDelimiter{\abs}{\lvert}{\rvert}
\DeclarePairedDelimiter{\norm}{\lVert}{\rVert}
\DeclarePairedDelimiter{\set}{\{}{\}}

% --- Hình và bảng -------------------------------------------------------------
\usepackage{graphicx}
\graphicspath{{figures/}}
\usepackage[font=small,labelfont=bf,format=hang]{caption}
\usepackage{subcaption}
\usepackage{booktabs}
\usepackage{array}
\usepackage{multirow}
\usepackage{tabularx}
\usepackage{longtable}
\usepackage{float}
\usepackage[table]{xcolor}
\usepackage{csvsimple}           % bảng đọc thẳng từ CSV (skill: Tables from CSV Data)
\usepackage{pgfplots}             % biểu đồ từ CSV (skill: charts-and-graphs)
\pgfplotsset{compat=1.18}

% --- Danh sách (nén khoảng cách) ----------------------------------------------
\usepackage{enumitem}
\setlist[itemize]{nosep, leftmargin=*, topsep=2pt, partopsep=0pt}
\setlist[enumerate]{nosep, leftmargin=*, topsep=2pt, partopsep=0pt}

% --- Giả mã -------------------------------------------------------------------
\usepackage{algorithm}
\usepackage{algpseudocode}
\floatname{algorithm}{Thuật toán}
\renewcommand{\listalgorithmname}{Danh sách thuật toán}
\algrenewcommand\algorithmicrequire{\textbf{Đầu vào:}}
\algrenewcommand\algorithmicensure{\textbf{Đầu ra:}}

% --- Đơn vị và số -------------------------------------------------------------
\usepackage{siunitx}
\sisetup{detect-all, output-decimal-marker={,}, per-mode=symbol}  % dấu phẩy thập phân theo văn bản tiếng Việt; đơn vị dạng kbit/s
\DeclareSIUnit{\dBm}{dBm}         % siunitx không có sẵn dBm
\DeclareSIUnit{\dB}{dB}

% --- Tài liệu tham khảo -------------------------------------------------------
\usepackage{csquotes}
\DeclareQuoteAlias{english}{vietnamese}   % csquotes chưa có kiểu ngoặc kép tiếng Việt
% biblatex chưa có chuỗi tiếng Việt ("vol.", "no.", "pp.") nên ánh xạ sang chuỗi tiếng Anh theo chuẩn IEEE
\usepackage[backend=biber, style=ieee, sorting=none, maxbibnames=6, autolang=none]{biblatex}
\DeclareLanguageMapping{vietnamese}{english}
% \addbibresource nằm trong main.tex để script compile_latex.sh của skill nhận ra cần chạy biber

% --- Liên kết và tham chiếu chéo (nạp gần cuối) -------------------------------
\usepackage[unicode, colorlinks=true, linkcolor=blue!60!black,
            citecolor=blue!60!black, urlcolor=blue!60!black]{hyperref}
\usepackage[nameinlink]{cleveref}
% cleveref chưa có tiếng Việt, đặt tên bằng tay
\crefname{figure}{Hình}{Hình}
\Crefname{figure}{Hình}{Hình}
\crefname{table}{Bảng}{Bảng}
\Crefname{table}{Bảng}{Bảng}
\crefname{equation}{công thức}{các công thức}
\Crefname{equation}{Công thức}{Các công thức}
\crefname{section}{Mục}{Mục}
\Crefname{section}{Mục}{Mục}
\crefname{subsection}{Mục}{Mục}
\Crefname{subsection}{Mục}{Mục}
\crefname{algorithm}{Thuật toán}{Thuật toán}
\Crefname{algorithm}{Thuật toán}{Thuật toán}
\crefname{theorem}{Định lý}{Định lý}
\Crefname{theorem}{Định lý}{Định lý}
\crefname{proposition}{Mệnh đề}{Mệnh đề}
\Crefname{proposition}{Mệnh đề}{Mệnh đề}
\crefname{definition}{Định nghĩa}{Định nghĩa}
\Crefname{definition}{Định nghĩa}{Định nghĩa}
\crefname{remark}{Nhận xét}{Nhận xét}
\Crefname{remark}{Nhận xét}{Nhận xét}
\crefname{lemma}{Bổ đề}{Bổ đề}
\Crefname{lemma}{Bổ đề}{Bổ đề}
\crefname{definition}{Định nghĩa}{Định nghĩa}
\Crefname{definition}{Định nghĩa}{Định nghĩa}
\crefname{proposition}{Mệnh đề}{Mệnh đề}
\Crefname{proposition}{Mệnh đề}{Mệnh đề}
